\documentclass[11pt]{article}
\usepackage[margin=1in]{geometry}
\usepackage{amsmath,amssymb,amsthm,mathtools}
\usepackage{graphicx}
\usepackage[hidelinks]{hyperref}
\usepackage{xcolor}

\newtheorem{theorem}{Theorem}
\newtheorem{corollary}[theorem]{Corollary}
\newtheorem{lemma}[theorem]{Lemma}
\newcommand{\R}{\mathbb R}
\newcommand{\TV}{\operatorname{TV}}

\title{A Smooth Counterexample to the Exponential\\
Total-Variation Conjecture for the Truncated Hilbert Transform}
\author{Automatic mathematics run \texttt{fa\_banach\_001}}
\date{9 August 2026}

\begin{document}
\setlength{\emergencystretch}{3em}
\maketitle

\begin{center}
\fbox{\parbox{0.91\textwidth}{
\textbf{Status.} Candidate counterexample, likely valid, for expert review.
The construction gives a full negative answer to equation (2.1) and to the
exponential stability conjecture after Theorem 2.1 of
Alaifari--Pierce--Steinerberger.  It works for every interval pair in their
Cases 3 and 4 and already uses real functions in $C_c^\infty(I)$.
}}
\end{center}

\section{Source conjecture}

For finite open intervals $I,J\subset\R$, the truncated Hilbert transform is
normalized here as
\[
 Hf(x)=\frac1\pi\operatorname{p.v.}\int_I\frac{f(y)}{x-y}\,dy,
 \qquad x\in J.
\]
The source is R. Alaifari, L. B. Pierce, and S. Steinerberger,
\emph{Lower bounds for the truncated Hilbert transform},
arXiv:1311.6845, published in \emph{Revista Matem\'atica Iberoamericana}
32 (2016), 23--56 \cite{source}.

On printed page 7, Section 2.2, for $I,J$ in Case 3 (a gap) or Case 4
(overlap), Theorem 2.1 establishes the
existence of some positive stability function for $BV(I)$ and then conjectures
an exponential lower rate.  Equation (2.1) states the corresponding concrete
inequality for weakly differentiable functions:
\begin{equation}\label{eq:conjecture}
 \|Hf\|_{L^2(J)}
 \ge c_1\exp\!\left(-c_2\frac{\|f'\|_{L^1(I)}}
                              {\|f\|_{L^2(I)}}\right)
       \|f\|_{L^2(I)},
\end{equation}
where $c_1,c_2>0$ may depend on $I,J$ but not on $f$.
The source evidence is reproduced below.

\begin{center}
\includegraphics[width=0.98\textwidth]{figures/open_problem_crop.png}
\end{center}

The paper itself notes in Sections 1.5 and 2.4 that Case 4 contains the gap
mechanism by restricting to functions compactly supported away from $J$.
That is precisely the localization used below.

\section{Proof intuition and result}

For $K>0$, define the maximal lower-stability envelope
\begin{equation}\label{eq:envelope}
 \mathcal S_{I,J}(K)=
 \inf\left\{
 \frac{\|Hf\|_{L^2(J)}}{\|f\|_{L^2(I)}}:
 0\ne f\in W^{1,1}(I),\quad
 \frac{\|f'\|_{L^1(I)}}{\|f\|_{L^2(I)}}\le K
 \right\}.
\end{equation}
The conjecture would imply
$\mathcal S_{I,J}(K)\ge c_1e^{-c_2K}$.

The counterexample is a train of disjoint smooth bumps with coefficients
\[
 1,-\binom m1,\binom m2,\ldots,(-1)^m.
\]
The $L^2$ norm of this train is governed by
$\sum_j\binom mj^2=\binom{2m}m$, whereas its total variation is governed by
$\sum_j\binom mj=2^m$.  Their quotient contributes only $m^{1/4}$.
Each bump has width comparable to $m^{-1}$, contributing another $m^{1/2}$,
so the normalized variation is $O(m^{3/4})$.

On the observation interval the Hilbert kernel is smooth.  The alternating
binomial coefficients are an $m$th finite difference, hence annihilate the
first $m-1$ orders of the kernel.  The remainder is geometrically small in
$m$.  Thus the data are $O(e^{-cm})$ while the regularity parameter is only
$O(m^{3/4})$, producing an $e^{-cK^{4/3}}$ upper obstruction.

\begin{theorem}[Smooth binomial-bump obstruction]\label{thm:main}
Let $I,J\subset\R$ be finite open intervals such that
$I\setminus\overline J$ contains a nonempty open interval.  This includes
every interval pair in Cases 3 and 4 of the source paper.  There are constants
$A,B,c>0$ and real functions $f_m\in C_c^\infty(I)$ such that, for all
sufficiently large integers $m$,
\begin{equation}\label{eq:kappa}
 \frac{\|f_m'\|_{L^1(I)}}{\|f_m\|_{L^2(I)}}
   \le A m^{3/4},
\end{equation}
and
\begin{equation}\label{eq:delta}
 \frac{\|Hf_m\|_{L^2(J)}}{\|f_m\|_{L^2(I)}}
   \le B e^{-cm}.
\end{equation}
Consequently, for suitable $B_0,c_0,K_0>0$,
\begin{equation}\label{eq:fourthirds}
 \mathcal S_{I,J}(K)\le B_0e^{-c_0K^{4/3}}
 \qquad(K\ge K_0).
\end{equation}
In particular, the conjectured inequality \emph{\eqref{eq:conjecture}} is
false, even on $C_c^\infty(I)$.
\end{theorem}

\section{Proof}

Choose a compact interval $K\Subset I\setminus\overline J$ with positive
length and put
\[
 D=\operatorname{dist}(K,J)>0.
\]
Choose $\rho>0$ so small that $2\rho<|K|$ and $\rho<D$.  Fix a point $a$
such that $[a,a+2\rho]\subset K$, and fix a nonzero real
$\phi\in C_c^\infty(0,1)$.  For $m\ge2$, put $h_m=\rho/m$ and define
\begin{equation}\label{eq:fm}
 f_m(y)=\sum_{j=0}^m(-1)^j\binom mj
 \phi\!\left(\frac{y-a-jh_m}{h_m}\right).
\end{equation}
All summands have disjoint supports, and their total span is
$(m+1)h_m<2\rho$.  Thus $f_m$ is real, nonzero, and belongs to
$C_c^\infty(K)\subset C_c^\infty(I)$.

\subsection*{Exact norm identities}

Disjointness and a change of variables give
\begin{equation}\label{eq:l2}
 \|f_m\|_{L^2(I)}^2
 =h_m\|\phi\|_2^2\sum_{j=0}^m\binom mj^2
 =h_m\|\phi\|_2^2\binom{2m}m,
\end{equation}
while
\begin{equation}\label{eq:l1}
 \|f_m'\|_{L^1(I)}
 =\|\phi'\|_1\sum_{j=0}^m\binom mj
 =2^m\|\phi'\|_1.
\end{equation}
Using the standard central-binomial estimate
\begin{equation}\label{eq:central}
 \binom{2m}m\ge\frac{4^m}{2\sqrt m}\qquad(m\ge1),
\end{equation}
we obtain
\begin{align*}
 \frac{\|f_m'\|_1}{\|f_m\|_2}
 &\le \sqrt2\,\frac{\|\phi'\|_1}{\|\phi\|_2}
       m^{1/4}h_m^{-1/2}\\
 &=\sqrt{\frac2\rho}\,\frac{\|\phi'\|_1}{\|\phi\|_2}
       m^{3/4}.
\end{align*}
This proves \eqref{eq:kappa}.

\subsection*{Finite-difference cancellation}

For $x\in J$, the support of $f_m$ is separated from $x$, so no principal
value is needed.  Changing variables in every summand of \eqref{eq:fm} gives
\begin{equation}\label{eq:kernel-sum}
 \pi Hf_m(x)=h_m\int_0^1\phi(t)
 \sum_{j=0}^m(-1)^j\binom mj
 \frac{dt}{x-a-h_m(j+t)}.
\end{equation}
For fixed $x,t$, let
\[
 q_{x,t}(s)=\frac1{x-a-h_m(s+t)}.
\]
The alternating sum in \eqref{eq:kernel-sum} is
$(-1)^m\Delta^m q_{x,t}(0)$.  Repeated use of the fundamental theorem of
calculus gives the exact identity
\begin{equation}\label{eq:difference}
 \Delta^m q(0)=\int_{[0,1]^m}
 q^{(m)}(u_1+\cdots+u_m)\,du_1\cdots du_m.
\end{equation}
Every point $a+h_m(s+t)$ occurring here lies in $K$.  Therefore
\begin{equation}\label{eq:derivative}
 |q_{x,t}^{(m)}(s)|
 =\frac{m!h_m^m}{|x-a-h_m(s+t)|^{m+1}}
 \le\frac{m!h_m^m}{D^{m+1}}.
\end{equation}
Combining \eqref{eq:kernel-sum}--\eqref{eq:derivative},
\begin{equation}\label{eq:pointwise}
 |Hf_m(x)|
 \le\frac{\|\phi\|_1}{\pi}
      \frac{m!h_m^{m+1}}{D^{m+1}}
 \qquad(x\in J).
\end{equation}

Taking the $L^2(J)$ norm and dividing by \eqref{eq:l2}, then using
\eqref{eq:central}, yields a constant $C=C(I,J,\phi)$ such that
\begin{equation}\label{eq:ratio-bound}
 \frac{\|Hf_m\|_{L^2(J)}}{\|f_m\|_{L^2(I)}}
 \le C\frac{m!h_m^{m+1/2}m^{1/4}}{2^mD^{m+1}}.
\end{equation}
Since $h_m=\rho/m$ and $m!\le m^m$, the right-hand side is at most
\begin{equation}\label{eq:geometric}
 C' m^{-1/4}\left(\frac{\rho}{2D}\right)^m.
\end{equation}
Our choice $\rho<D$ makes the base smaller than $1/2$, proving
\eqref{eq:delta}.

Finally let $A$ be the constant in \eqref{eq:kappa}.  For sufficiently large
$K$, choose an integer
\[
 m=\left\lfloor\left(\frac{K}{2A}\right)^{4/3}\right\rfloor.
\]
Then $Am^{3/4}\le K$ and $m\ge c'K^{4/3}$ after increasing the lower
threshold.  The definition \eqref{eq:envelope} and \eqref{eq:delta} give
\eqref{eq:fourthirds}.

If \eqref{eq:conjecture} held, then the same functions would satisfy
\[
 \frac{\|Hf_m\|_2}{\|f_m\|_2}
 \ge c_1e^{-c_2Am^{3/4}},
\]
contradicting the geometric upper bound \eqref{eq:geometric}.  This completes
the proof. \qed

\section{Consequences and quantitative window}

\begin{corollary}[Negative answer to both source formulations]
For every Case 3 or Case 4 interval pair, the exponential conjecture after
Theorem 2.1 and its weakly differentiable formulation (2.1) are false.  The
failure is not caused by jumps or roughness: it occurs inside
$C_c^\infty(I)$.
\end{corollary}

For $f\in W^{1,1}(I)$, total variation equals $\|f'\|_1$, so the theorem
directly applies to the $BV$ stability parameter used in Theorem 2.1.  On the
other hand, Theorem 2.4 of the source proves a lower bound of the form
\[
 \mathcal S_{I,J}(K)\ge c_1e^{-c_2K^2}.
\]
Together with \eqref{eq:fourthirds}, the presently justified window is
\[
 c_1e^{-c_2K^2}
 \ \le\ \mathcal S_{I,J}(K)
 \ \le\ B_0e^{-c_0K^{4/3}}
 \qquad(K\gg1).
\]
Determining the sharp exponent or non-power correction inside this window
remains open.  The present result is therefore a full refutation and a
quantitative obstruction, not a full characterization of the optimal modulus.

\section{Verification and novelty bounds}

The accompanying script \path{code/verify_binomial_bumps.py} verifies the
two binomial identities and the reciprocal-kernel finite-difference identity
using exact integer/rational arithmetic.  It also checks the $m^{3/4}$ scaling
and geometric decay numerically.  The computation is a sanity check and is
not used in the proof.

On 9 August 2026, the four local run indexes were searched for the arXiv id,
exact title, displayed conjecture, and core keyword combinations.  Exact-title
and keyword searches of arXiv and the web found the source, related truncated
transform papers, and the dyadic model of Jaming--Pozzi--Wick \cite{dyadic},
but no refutation, alternating-binomial construction, or $4/3$ obstruction.
The dyadic work concerns a different operator and does not answer the present
continuous conjecture.  This is a bounded novelty check, not certification;
MathSciNet, zbMATH, full citation graphs, and author correspondence were not
exhaustively checked.

\section{Human-review recommendation}

\textbf{Recommendation: promote after expert confirmation.}  The main review
points are: (i) the source's Case 4 orientation and the separated subinterval
of $I\setminus\overline J$; (ii) the signs and derivative order in
\eqref{eq:difference}--\eqref{eq:derivative}; (iii) the normalization in
\eqref{eq:ratio-bound}; and (iv) whether the construction has appeared outside
the bounded novelty search.  None of these checks revealed a defect in the
present audit.

\begin{thebibliography}{9}
\bibitem{source}
R. Alaifari, L. B. Pierce, and S. Steinerberger,
\emph{Lower bounds for the truncated Hilbert transform},
Rev. Mat. Iberoam. 32 (2016), no. 1, 23--56;
arXiv:1311.6845; doi:10.4171/RMI/880.

\bibitem{dyadic}
P. Jaming, E. Pozzi, and B. D. Wick,
\emph{Lower bounds for the dyadic Hilbert transform},
Ann. Fac. Sci. Toulouse Math. (6) 26 (2017), no. 1, 185--204;
arXiv:1605.05511.
\end{thebibliography}

\end{document}
